% ============ 7. Referencias (APA 7) ============
\refapa{Cañibano, C., \& Bozeman, B. (2009). Curriculum vitae method in science policy and research evaluation: The state-of-the-art. \emph{Research Policy, 38}(2), 259--266. https://doi.org/10.1016/j.respol.2008.09.014}

\refapa{Categorización de grupos e investigadores ante el sistema nacional de ciencia, tecnología e innovación en Colombia. (2023). \emph{Universidad y Sociedad, 15}(5). http://scielo.sld.cu/scielo.php?pid=S2218-36202023000500133}

\refapa{Clauset, A., Arbesman, S., \& Larremore, D. B. (2015). Systematic inequality and hierarchy in faculty hiring networks. \emph{Science Advances, 1}(1), e1400005. https://doi.org/10.1126/sciadv.1400005}

\refapa{CoARA. (2022). \emph{Agreement on reforming research assessment}. https://coara.eu/agreement/the-agreement-full-text/}

\refapa{Crenshaw, K. (1989). Demarginalizing the intersection of race and sex: A Black feminist critique of antidiscrimination doctrine, feminist theory and antiracist politics. \emph{University of Chicago Legal Forum, 1989}, 139--167.}

\refapa{Departamento Nacional de Planeación. (2021). \emph{CONPES 4069: Política Nacional de Ciencia Abierta}. https://colaboracion.dnp.gov.co/CDT/Conpes/Económicos/4069.pdf}

\refapa{Deville, P., Song, C., Sinatra, R., Wang, D., \& Barabási, A.-L. (2014). Career on the move: Geography, stratification, and scientific impact. \emph{Proceedings of the National Academy of Sciences, 111}(42), 14472--14476. https://doi.org/10.1073/pnas.1407385111}

\refapa{DORA. (2012). \emph{San Francisco Declaration on Research Assessment}. https://sfdora.org/}

\refapa{Fortunato, S., Bergstrom, C. T., Börner, K., Evans, J. A., Helbing, D., Milojević, S., Petersen, A. M., Radicchi, F., Sinatra, R., Uzzi, B., et al. (2018). Science of science. \emph{Science, 359}(6379), eaao0185. https://doi.org/10.1126/science.aao0185}

\refapa{Frenken, K., Hardeman, S., \& Hoekman, J. (2009). Spatial scientometrics: Towards a cumulative research program. \emph{Journal of Informetrics, 3}(3), 222--232. https://doi.org/10.1016/j.joi.2009.03.005}

\refapa{Hicks, D., Wouters, P., Waltman, L., de Rijcke, S., \& Rafols, I. (2015). The Leiden Manifesto for research metrics. \emph{Nature, 520}(7548), 429--431. https://doi.org/10.1038/520429a}

\refapa{Kemeny, J. G., \& Snell, J. L. (1976). \emph{Finite Markov chains}. Springer.}

\refapa{Larivière, V., Ni, C., Gingras, Y., Cronin, B., \& Sugimoto, C. R. (2013). Bibliometrics: Global gender disparities in science. \emph{Nature, 504}(7479), 211--213. https://doi.org/10.1038/504211a}

\refapa{Ley 1581 de 2012. Por la cual se dictan disposiciones generales para la protección de datos personales. Congreso de la República de Colombia.}

\refapa{Mena-Chalco, J. P., \& Cesar Júnior, R. M. (2009). ScriptLattes: integração de dados curriculares à plataforma Lattes. \emph{Conferência sobre Tecnologia, Cultura e Sociedade}.}

\refapa{Ministerio de Ciencia, Tecnología e Innovación. (2021). \emph{Investigadores reconocidos por convocatoria} [Conjunto de datos]. datos.gov.co. https://www.datos.gov.co/resource/bqtm-4y2h.csv}

\refapa{Ministerio de Ciencia, Tecnología e Innovación. (2021). \emph{Producción de grupos de investigación} [Conjunto de datos]. datos.gov.co. https://www.datos.gov.co/resource/33dq-ab5a.csv}

\refapa{Newman, M. E. J. (2001). The structure of scientific collaboration networks. \emph{Proceedings of the National Academy of Sciences, 98}(2), 404--409. https://doi.org/10.1073/pnas.98.2.404}

\refapa{Petersen, A. M., Riccaboni, M., Stanley, H. E., \& Pammolli, F. (2012). Persistence and uncertainty in the academic career. \emph{Proceedings of the National Academy of Sciences, 109}(14), 5213--5218. https://doi.org/10.1073/pnas.1121429109}

\refapa{Revista Iberoamericana de Ciencia, Tecnología y Sociedad. (2010). Dos países latinoamericanos cuentan actualmente con sistemas de información curricular consolidados: Brasil y Colombia. \emph{CTS, 5}(13). https://www.revistacts.net/wp-content/uploads/2020/01/vol5-nro13-cts13.pdf}

\refapa{Rocher, L., Hendrickx, J. M., \& de Montjoye, Y.-A. (2019). Estimating the success of re-identifications in incomplete datasets: Generating insights from attacks. \emph{Nature Communications, 10}, 3069. https://doi.org/10.1038/s41467-019-10933-3}

\refapa{Sinatra, R., Wang, D., Deville, P., Song, C., \& Barabási, A.-L. (2016). Quantifying the evolution of individual scientific impact. \emph{Science, 354}(6312), aaf5239. https://doi.org/10.1126/science.aaf5239}

\refapa{Souza, R., Skluzacek, T., Wilkinson, S., Ziatdinov, M., \& da Silva, R. (2023). \emph{Towards lightweight data integration using multi-workflow provenance and data observability}. IEEE eScience (preprint). https://arxiv.org/abs/2308.09004}

\refapa{Sugimoto, C. R., Robinson-García, N., Murray, D. S., Yegros-Yegros, A., Costas, R., \& Larivière, V. (2017). Scientists have most impact when they're free to move. \emph{Nature, 550}(7674), 29--31. https://doi.org/10.1038/550029a}

\refapa{Theil, H. (1967). \emph{Economics and information theory}. North-Holland.}

\refapa{Wang, D., \& Barabási, A.-L. (2021). \emph{The science of science}. Cambridge University Press.}

\refapa{Wang, R. Y., \& Strong, D. M. (1996). Beyond accuracy: What data quality means to data consumers. \emph{Journal of Management Information Systems, 12}(4), 5--33. https://doi.org/10.1080/07421222.1996.11518099}

\refapa{Wilkinson, M. D., Dumontier, M., Aalbersberg, I. J., Appleton, G., Axton, M., Baak, A., et al. (2016). The FAIR Guiding Principles for scientific data management and stewardship. \emph{Scientific Data, 3}, 160018. https://doi.org/10.1038/sdata.2016.18}
